\section{Prior Concluding Discussion (R10)}
Neural Bellman Operators organize computation around policy evaluation, feasible improvement, and independent economic verification. Recursive utility, endogenous preferences, sophisticated temporal selves, and dynamic strategic interaction remain part of that economic formulation. The representation is neural, but an economic claim concerns the policy actually deployed and the error that has been bounded.

The guarded map specifies how a nodal evaluation tolerance is enforced and records the indexed corrections it requires. Signed derivative enclosures reduce complete-action verification work without dropping unresolved boxes or differentiating across a stopping discontinuity. Their numerical contribution is identified using identical frozen policies and matched one-step comparisons. Fresh training and recertification on genuinely nested grids are kept distinct from the full-domain failure of interpolated coarse policies. The compensation account translates utility accuracy into a stated financed-transfer experiment rather than a misleading common percentage.

The dense capital experiment now has a feasible-policy lower benchmark and a continuous-time optimal upper benchmark on the original state space. Those historical absolute brackets are distinct from the new continuous-time neural-policy bound in Theorem~\ref{thm:r10tube}. The latter uses the original capital economy and full adapted comparison class, with a schedule-centred actor architecture. Paired simulations remain discretized policy comparisons. Likewise, the complete-action preference certificates do not yet bound all continuous-state, continuous-time, and stopping-monitoring defects. The evidence does not establish that actors outperform the strongest low-dimensional search and projection alternatives; those comparisons remain in the record. The contribution is a specified neural Bellman computation with verified evaluation and improvement outputs, measured costs, and separate economic domains for each conclusion.

